% GENERATED FILE. Edit canonical lessons, then run npm run generate:books.
% canonical-source-hash: fnv1a64:ad0c56c6ef3a10be
% canonical-lessons: GU-C363-akdo, GU-C363-gantari, GU-C363-taran, GU-C363-sankhya, GU-C363-ghatadvu

\chapter{What the Figures Say}
\label{ch:gu-what-the-figures-say}
\hlchaptermodality{363}

\begin{chapteropening}
\textbf{By the end of this chapter:} \emph{I can say what the numbers show.}

Complete the last lesson of chapter 363: I can say what the numbers show.
\end{chapteropening}

\section[\texorpdfstring{\gu{આંકડો} (ā̃kḍo)}{આંકડો (ā̃kḍo)}]{\gu{આંકડો} (ā̃kḍo) — a figure (a number)}

\label{lesson:GU-C363-akdo}



\begin{quote}
Before the new one: say the Gujarati for a sample, then the Gujarati for to deduce.
\end{quote}

\subsection*{You'll want to know: \gu{આંકડો}}
\textbf{\gu{આંકડો}} — \emph{ā̃kḍo} — \textquotedblleft{}a figure (a number)\textquotedblright{}.

\begin{quote}
\gu{આ} \gu{આંકડા} \gu{દર્શાવે} \gu{છે} \gu{કે} \gu{ગામમાં} \gu{વરસાદ} \gu{ઓછો} \gu{થયો}.
\end{quote}

\begin{quote}
\emph{ā ā̃kḍā darśāve che ke gāmmā̃ varsād ocho thayo.}
\end{quote}

\begin{quote}
These figures show that the village has had less rain.
\end{quote}

A masculine noun.

\begin{cousinweb}
\textbf{\gu{આંકડો}} (\emph{ā̃kḍo}) is built on Sanskrit \emph{aṅka}, \textquotedblleft{}a mark, a numeral\textquotedblright{} (Monier-Williams).
\end{cousinweb}

\begin{grammarlens}[title={What you've built}]
One more word for giving reasons and opinions.
\end{grammarlens}

\subsection*{Your turn}
\begin{itemize}
\raggedright
  \item \emph{Say it:} \emph{ā̃kḍo}
  \item \emph{Say it:} \emph{ā̃kḍo}, once more
  \item \emph{Say it:} \emph{ā ā̃kḍā darśāve che ke gāmmā̃ varsād ocho thayo.} — These figures show that the village has had less rain.
  \item \emph{Recall:} say the Gujarati for a sample, then the Gujarati for to deduce, then say \emph{ā̃kḍo} again
\end{itemize}

\begin{tcolorbox}[breakable,colback=teal!4,colframe=teal!35!black,title={Before you move on}]
What does \gu{આંકડો} mean? (\textquotedblleft{}A figure (a number)\textquotedblright{}.) Say it once more.
\end{tcolorbox}
\section[\texorpdfstring{\gu{ગણતરી} (gaṇtarī)}{ગણતરી (gaṇtarī)}]{\gu{ગણતરી} (gaṇtarī) — a calculation}

\label{lesson:GU-C363-gantari}



\begin{quote}
Before the new one: say the Gujarati for to deduce, then the Gujarati for a figure.
\end{quote}

\subsection*{You'll want to know: \gu{ગણતરી}}
\textbf{\gu{ગણતરી}} — \emph{gaṇtarī} — \textquotedblleft{}a calculation\textquotedblright{}.

\begin{quote}
\gu{મારી} \gu{ગણતરી} \gu{સાચી} \gu{નીકળી}.
\end{quote}

\begin{quote}
\emph{mārī gaṇtarī sācī nīkḷī.}
\end{quote}

\begin{quote}
My calculation turned out right.
\end{quote}

A feminine noun.

\begin{cousinweb}
\textbf{\gu{ગણતરી}} (\emph{gaṇtarī}) is built on \textbf{\gu{ગણવું}} (\emph{gaṇvũ}), \textquotedblleft{}to count\textquotedblright{}, from Sanskrit \emph{gaṇayati}, \textquotedblleft{}counts\textquotedblright{} (CDIAL 3993).
\end{cousinweb}

\begin{grammarlens}[title={What you've built}]
One more word for giving reasons and opinions.
\end{grammarlens}

\subsection*{Your turn}
\begin{itemize}
\raggedright
  \item \emph{Say it:} \emph{gaṇtarī}
  \item \emph{Say it:} \emph{gaṇtarī}, once more
  \item \emph{Say it:} \emph{mārī gaṇtarī sācī nīkḷī.} — My calculation turned out right.
  \item \emph{Recall:} say the Gujarati for to deduce, then the Gujarati for a figure, then say \emph{gaṇtarī} again
\end{itemize}

\begin{tcolorbox}[breakable,colback=teal!4,colframe=teal!35!black,title={Before you move on}]
What does \gu{ગણતરી} mean? (\textquotedblleft{}A calculation\textquotedblright{}.) Say it once more.
\end{tcolorbox}
\section[\texorpdfstring{\gu{તારણ} (tāraṇ)}{તારણ (tāraṇ)}]{\gu{તારણ} (tāraṇ) — a finding (a conclusion)}

\label{lesson:GU-C363-taran}



\begin{quote}
Before the new one: say the Gujarati for a figure, then the Gujarati for a calculation.
\end{quote}

\subsection*{You'll want to know: \gu{તારણ}}
\textbf{\gu{તારણ}} — \emph{tāraṇ} — \textquotedblleft{}a finding (a conclusion)\textquotedblright{}.

\begin{quote}
\gu{આ} \gu{સંશોધનનું} \gu{તારણ} \gu{શું} \gu{છે}?
\end{quote}

\begin{quote}
\emph{ā sanśodhannũ tāraṇ śũ che?}
\end{quote}

\begin{quote}
What did this research find?
\end{quote}

A neuter noun.

\begin{cousinweb}
\textbf{\gu{તારણ}} (\emph{tāraṇ}) CDIAL lists Gujarati \emph{tāraṇ}, \textquotedblleft{}a taking across, a pledge\textquotedblright{}, under Sanskrit \emph{tāraṇa}, \textquotedblleft{}crossing, safe passage\textquotedblright{} (CDIAL 5795); \textquotedblleft{}a conclusion\textquotedblright{}, what is drawn out of the facts, is a sense CDIAL does not give.
\end{cousinweb}

\begin{grammarlens}[title={What you've built}]
One more word for giving reasons and opinions.
\end{grammarlens}

\subsection*{Your turn}
\begin{itemize}
\raggedright
  \item \emph{Say it:} \emph{tāraṇ}
  \item \emph{Say it:} \emph{tāraṇ}, once more
  \item \emph{Say it:} \emph{ā sanśodhannũ tāraṇ śũ che?} — What did this research find?
  \item \emph{Recall:} say the Gujarati for a figure, then the Gujarati for a calculation, then say \emph{tāraṇ} again
\end{itemize}

\begin{tcolorbox}[breakable,colback=teal!4,colframe=teal!35!black,title={Before you move on}]
What does \gu{તારણ} mean? (\textquotedblleft{}A finding (a conclusion)\textquotedblright{}.) Say it once more.
\end{tcolorbox}
\section[\texorpdfstring{\gu{સંખ્યા} (sankhyā)}{સંખ્યા (sankhyā)}]{\gu{સંખ્યા} (sankhyā) — a count; a number of people or things}

\label{lesson:GU-C363-sankhya}



\begin{quote}
Before the new one: say the Gujarati for a calculation, then the Gujarati for a finding.
\end{quote}

\subsection*{You'll want to know: \gu{સંખ્યા}}
\textbf{\gu{સંખ્યા}} — \emph{sankhyā} — \textquotedblleft{}a count; a number of people or things\textquotedblright{}.

\begin{quote}
\gu{આ} \gu{વરસે} \gu{વિદ્યાર્થીઓની} \gu{સંખ્યા} \gu{વધી}.
\end{quote}

\begin{quote}
\emph{ā varse vidyārthīonī sankhyā vadhī.}
\end{quote}

\begin{quote}
The number of students went up this year.
\end{quote}

A feminine noun.

\begin{cousinweb}
\textbf{\gu{સંખ્યા}} (\emph{sankhyā}) is a tatsama, borrowed from Sanskrit \emph{saṃkhyā}, \textquotedblleft{}reckoning, number\textquotedblright{} (Monier-Williams).
\end{cousinweb}

\begin{grammarlens}[title={What you've built}]
One more word for giving reasons and opinions.
\end{grammarlens}

\subsection*{Your turn}
\begin{itemize}
\raggedright
  \item \emph{Say it:} \emph{sankhyā}
  \item \emph{Say it:} \emph{sankhyā}, once more
  \item \emph{Say it:} \emph{ā varse vidyārthīonī sankhyā vadhī.} — The number of students went up this year.
  \item \emph{Recall:} say the Gujarati for a calculation, then the Gujarati for a finding, then say \emph{sankhyā} again
\end{itemize}

\begin{tcolorbox}[breakable,colback=teal!4,colframe=teal!35!black,title={Before you move on}]
What does \gu{સંખ્યા} mean? (\textquotedblleft{}A count; a number of people or things\textquotedblright{}.) Say it once more.
\end{tcolorbox}
\section[\texorpdfstring{\gu{ઘટાડવું} (ghaṭāḍvũ)}{ઘટાડવું (ghaṭāḍvũ)}]{\gu{ઘટાડવું} (ghaṭāḍvũ) — to bring down}

\label{lesson:GU-C363-ghatadvu}



\begin{quote}
Before the new one: say the Gujarati for a finding, then the Gujarati for a count.
\end{quote}

\subsection*{You'll want to know: \gu{ઘટાડવું}}
\textbf{\gu{ઘટાડવું}} — \emph{ghaṭāḍvũ} — \textquotedblleft{}to bring down\textquotedblright{}.

\begin{quote}
\gu{સરકાર} \gu{પેટ્રોલનો} \gu{ભાવ} \gu{ઘટાડશે}.
\end{quote}

\begin{quote}
\emph{sarkār peṭrolno bhāv ghaṭāḍśe.}
\end{quote}

\begin{quote}
The government will bring down the price of petrol.
\end{quote}

Past, with \gu{તેણે} (he or she …): \gu{તેણે} \gu{ભાવ} \gu{ઘટાડ્યો}, \emph{teṇe bhāv ghaṭāḍyo}.

\begin{cousinweb}
\textbf{\gu{ઘટાડવું}} (\emph{ghaṭāḍvũ}) is the causative of \textbf{\gu{ઘટવું}} (\emph{ghaṭvũ}), \textquotedblleft{}to go down\textquotedblright{}, from a reconstructed \emph{ghaṭṭati}, \textquotedblleft{}decreases\textquotedblright{} (CDIAL 4415).
\end{cousinweb}

\begin{grammarlens}[title={What you've built}]
One more word for giving reasons and opinions.
\end{grammarlens}

\subsection*{Your turn}
\begin{itemize}
\raggedright
  \item \emph{Say it:} \emph{ghaṭāḍvũ}
  \item \emph{Say it:} \emph{ghaṭāḍvũ}, once more
  \item \emph{Say it:} \emph{sarkār peṭrolno bhāv ghaṭāḍśe.} — The government will bring down the price of petrol.
  \item \emph{Recall:} say the Gujarati for a finding, then the Gujarati for a count, then say \emph{ghaṭāḍvũ} again
\end{itemize}

\begin{tcolorbox}[breakable,colback=teal!4,colframe=teal!35!black,title={Before you move on}]
What does \gu{ઘટાડવું} mean? (\textquotedblleft{}To bring down\textquotedblright{}.) Say it once more.
\end{tcolorbox}
